\section{Related Work}
\label{sec:related_work}

\paragraph{Single-instance counting and tempo modeling.}
Early learned counters estimated a local cycle length in an online
window~\cite{levy2015live}, while motion-centric methods used time-varying
flow differentials and wavelets to address non-stationary
repetition~\cite{runia2018realworld}. RepNet used temporal self-similarity as an
intermediate representation bottleneck and introduced Countix as an in-the-wild
benchmark for class-agnostic repetition counting~\cite{dwibedi2020repnet}.
TransRAC regresses a density map from
multi-scale temporal correlations~\cite{hu2022transrac}, and DeTRC represents
cycles with dynamic action queries under a single-primary-action assumption
~\cite{ma2026detrc}. ESCounts uses exemplar-conditioned correspondence during
training and learns a latent used for exemplar-free, zero-shot inference
~\cite{sinha2024escounts}; SkimFocusNet combines
coarse contextual guidance with fine-grained views. Multi-RepCount evaluates
exemplar-specified counting of one target action category in videos containing
multiple repetitive-action categories, whereas MultiRep evaluates per-person
counting for simultaneous human instances
~\cite{zhao2025skimfocus,tang2024multicounter}.

\paragraph{Pose-based and label-efficient counting.}
Pose representations reduce appearance dependence and expose articulated
motion directly. PoseRAC models actions through salient poses, uses foundation
generative models to synthesize pose-conditioned training data, and
incorporates CLIP text embeddings for open-set counting, enabling zero-shot
prediction without real benchmark training samples~\cite{yao2025poserac};
zero-shot prediction should not be conflated with a self-supervised training
objective. D$^2$-STX instead fuses RGB and pose through
decoupled spatial and temporal cross-attention~\cite{wang2026d2stx}. PAMS
learns from unlabeled skeleton sequences using period-adaptive multi-scale
temporal consistency and combines inference-time counting experts
~\cite{gao2026pams}. Its expert consensus is not a learned fast--medium--slow
mixture, and its single-stream formulation does not preserve multiple
identities. Our intended use of pose and PAMS-TCC is therefore a foundation,
not a novelty claim; the proposed distinction is track-indexed local routing in
MRAC under a no-count-label objective. Local routing and normalized overlap-add
are generic operations and are not claimed as independent inventions; any
defensible novelty must come from their audited role in the stated MRAC
supervision and non-stationary-tempo setting.

\paragraph{Multi-instance repetition counting.}
MultiCounter defines MRAC and introduces MultiRep, learning detection,
association, period localization, and person-wise counting jointly
~\cite{tang2024multicounter}. MultiCounter+ improves robustness and efficiency
with a dynamically updated Siamese query and anchor-guided similarity
reconstruction that represents long and short periods
~\cite{tang2026multicounterplus}. These systems already address asynchronous
people, variable speeds, pauses, and identity continuity. Our proposed setting
differs along a narrower axis: an upstream tracker supplies masked pose tracks,
the counting objective is intended to avoid person-wise count and period
labels, and tempo specialization varies across local windows of each identity.
This modular design is not described as end-to-end, and any comparison to
native MRAC systems must disclose the different input and supervision
protocols. Appendix~\ref{tab:prior_scope} makes these boundaries explicit; it
also prevents generic local routing or overlap-add from being presented as a
standalone novelty.
